\subsection{Metas del proyecto}
La meta central del proyecto es realizar una exploración
sistemática de un problema de optimización de alta
dimensionalidad (derivado de un modelo de teoría de cuerdas
tipo IIB) mediante la implementación y comparación de diversos
métodos inspirados en la naturaleza (NIM, por sus siglas en
inglés). Esta aproximación permite abordar el problema desde
una perspectiva computacional, evaluando qué familias de
algoritmos se adaptan mejor a la geometría particular de una
función potencial proveniente de la K-teoría.

El proyecto contempla la implementación de un conjunto de NIM
que incluye algoritmos genéticos, optimización por enjambre de
partículas, algoritmos de estimación de distribución, colonia
de hormigas, entre otros. Cada método será evaluado sobre las
mismas bases de datos de configuraciones iniciales para
garantizar la comparabilidad de los resultados, utilizando
métricas de eficiencia computacional, calidad de las
soluciones encontradas y capacidad para identificar nuevos
vacíos estables no reportados previamente en la literatura.

En cuanto a la formación de recursos humanos, el proyecto
contribuye al desarrollo de capacidades en modelado matemático
avanzado, optimización computacional, manejo de herramientas
de cálculo simbólico y programación científica. Estas
competencias, al ser transferibles a otras áreas de la ciencia
y la ingeniería, inciden en la consolidación de capital humano
especializado en métodos computacionales aplicados a problemas
de ciencia básica y de frontera.

Como producto tangible, se espera generar código científico
(liberado como software de acceso abierto) que implementa la
función objetivo y los algoritmos de optimización,
constituyendo una herramienta reusable para la comunidad
científica nacional e internacional. Asimismo, se buscará
publicar los resultados en una revista indexada en el Sistema
de Clasificación de Revistas Mexicanas de Ciencia y Tecnología
de CONAHCYT, contribuyendo de esta manera a la difusión del
conocimiento generado en el país.
